\documentclass[10pt,a4paper]{amsart}
\usepackage[margin=19mm]{geometry}
\usepackage{amsmath,amssymb,mathtools}
\usepackage[scr=rsfs,cal=euler]{mathalfa}
\usepackage{xfrac}
\usepackage[unicode,hidelinks,bookmarks=false]{hyperref}
\newcommand{\ie}{i.e.\ }
\newcommand{\cf}{cf.\ }
\newcommand{\resp}{respectively\ }
\newcommand{\Subsection}{\subsection}
\providecommand{\nobreakdash}{\nobreak\hbox{-}}
\setlength{\emergencystretch}{2em}
\multlinegap0pt
\numberwithin{equation}{section}
\newtheorem{theorem}{Theorem}[section]
\allowdisplaybreaks
\title[Analysis of an HIV infection model with delays: original Theorem 3.4]{Analysis of an HIV infection model with delays: the original Theorem 3.4 and its complete original proof}
\author{Guangsheng Lai, Zhiting Xu}
\date{Source: Electronic Journal of Differential Equations 2026, No. 20, pp. 1--14; DOI 10.58997/ejde.2026.20; publisher version}
\begin{document}
\maketitle
\noindent\emph{Electronic Journal of Differential Equations}, Vol. 2026 (2026), No. 20, pp. 1--14.\newline
ISSN: 1072-6691. DOI \href{https://doi.org/10.58997/ejde.2026.20}{10.58997/ejde.2026.20}. This work is licensed under a CC BY 4.0 license.\par\smallskip
\noindent\textit{Editorial scope.} Counted claim: exactly one, the article's original Theorem 3.4 as printed (source0.tex lines 674--678, the single primary statement), the paper's main global-stability result for the delayed HIV infection model with CTL immune response and immune impairment, together with its complete original proof (source0.tex lines 680--835, the single primary proof). Everything the proof uses is retained verbatim and uncounted: the whole of the original Section 1, Introduction (55--199), with the three successively refined models, the delayed model (1.4) itself and the description of every parameter and delay; the whole of the original Section 2, Preliminaries (200--448), that is the well-posedness Theorem 2.1 with its complete original proof, the construction of the basic reproduction number from the next infection operator, the existence and uniqueness of the equilibria and Theorem 2.2; and the whole of the original Section 3 up to the counted theorem (449--673), that is the characteristic-equation Theorems 3.1 and 3.2 with their complete original proofs, the global-stability Theorem 3.3 with its complete original proof, and the fundamental function of the Lyapunov functionals. Because no content line of the original Sections 1 to 3 is dropped, every printed equation, theorem and section number of the delivered text coincides with the published version: the counted theorem is printed as Theorem 3.4, the support theorems as Theorem 2.1, Theorem 2.2, Theorem 3.1, Theorem 3.2 and Theorem 3.3, and the model as (1.4). Excluded from this unit: the whole of the original Section 4, Numerical simulations (837--904), which carries the parameter table and the two illustrations, and the whole of the original Section 5, Summary (905--934), as well as the acknowledgements and the printed bibliography, whose entries are transcribed into the delivery in published order with their published numbers. No figure environment, no image inclusion line, no table and no graphic asset is delivered, and none is redistributed. Printed slips are kept literal and no mathematics has been written, repaired or re-derived: in Theorem 2.2 the infection-free equilibrium is said to be defined by (2.1) although the display that defines it carries the number (2.4) (line 438); in the same theorem the two coordinates of the infection equilibrium are given by the same printed formula and the symbol printed there is u where the model uses mu (line 444, repeated in the proof of Theorem 3.2 at line 549); in the display for the derivative of the first half of the counted Lyapunov functional the delayed infection term is printed with $\tau$ in place of $\tau_1$ (line 703); the coefficient of the last term in that display carries the factor (1 - x*/x(t)) while the expansion that follows writes the same contribution as (z(t) - z*)/z(t) (lines 708 and 719--720); the final expression prints the product of the immune-impairment factor with a doubled y, as omega y* y z(t) (lines 811--812 and 822--823); and one logarithm in the later display is printed as italic ln instead of the upright operator used everywhere else (line 787). In the support material the cancellation of the two reproduction-number terms is printed with $\tau_2$ in one line and $\tau_1$ in the next (lines 642 and 646); it is recorded as printed and not repaired. One source citation command is printed with a space between the control sequence and its argument (line 152); the editorial note that follows lists the five works it cites, so that the delivered reference list is exactly the set of works cited by the excerpt.
\noindent\textit{Editorial note on one source citation.} The source prints one citation command with a space between the control sequence and its argument (source0.tex line 152, inside the sentence on the logistic proliferation of the target cells). So that the delivered reference list is exactly the set of works cited by the excerpt, the five works carried by that command are listed here: \cite{LGM,LS,PC,WLL,ZZZ}.\par\smallskip
\section{Introduction}

Acquired immune deficiency syndrome (AIDS)
is a devastating infectious disease caused by infection with a virus that attacks 
the body's immune system. It takes the most important CD$4^{+}$T lymphocytes in 
the human immune system as the main target and destroys a large number  
of cells, so that the human body losses its immune function. 
Therefore, human infection leads to immune deficiency and being prone to
various diseases, and the occurrence of malignant tumors.
A classical viral infection model was proposed in \cite{BM},
\begin{equation}\label{1.1}
\begin{gathered}
\dot{x}(t)=\lambda-dx(t)-\beta x(t)v(t),  \\
\dot{y}(t)=\beta x(t)v(t)-ay(t),\\
\dot{v}(t)=ky(t)-\mu v(t),
\end{gathered}
\end{equation}
where $x(t)$, $y(t)$, and $v(t)$ represent the concentrations of susceptible cells, 
infected cells, and free HIV virions, respectively.
The constants $\lambda$ and $d$  denote the production  and mortality rates of 
susceptible cells; $\beta$ is the effective
contact rate between susceptible cells and HIV virions; 
$a$ denotes the mortality rate of infected target cells; 
$k$ represents the average rate at which each infected target
cell produces free virions; $\mu$ is the virus clearance rate.  
By constructing Lyapunov functions, Korobeinikov \cite{K}  established  the global 
dynamics of  model \eqref{1.1}.

To  study the pathogenesis and transmission mechanism of AIDS, many researchers 
have studied the dynamic model of HIV infection in the host
(see \cite{LW2, NA, NB, Sa, WL} and the references cited therein).
 We note that \cite{NA} focus on eliminating or controlling the spread of the virus, 
taking into account the immune response that the virus has during
infection. The rapid and nonspecific immune response initially induced in the host 
during viral infection is mainly achieved by natural killer cells. 
Natural killer cells are part of the immune system,
and they play an important role in the early stages of viral infection. 
It is able to quickly identify and attack cells infected by the virus, thus 
preventing its  replication and spread. 
But in most viral infections, cytotoxic T lymphocyte (CTL) cells that attack infected 
cells and antibody cells that attack the virus play a  key role in antiviral defense. 
To investigate the role of the population dynamics of viral infection with 
CTL response, Nowak and Bangham \cite{NB} introduced a  mathematical model describing 
the basic dynamics of the interaction between activated CD$4^{+}$ T cells $x(t)$, 
infected CD$4^{+}$ T cells $y(t)$, viruses  $v(t)$ and
 immune cells $z(t)$. The model is 
\begin{equation}\label{1.2}
\begin{gathered}
\dot{x}(t)=\lambda-dx(t)-\beta x(t)v(t),\\
\dot{y}(t)=\beta x(t)v(t)-ay(t)-py(t)z(t),\\
\dot{v}(t)=ky(t)-\mu v(t),\\
\dot{z}(t)=cy(t)-bz(t).
\end{gathered}
\end{equation}
On the basis of this model, Song and Neumann \cite{SN} proposed the 
assumption that the clearance rate of each cell to the infected target cell for the 
cellular immune process is constant, that is, the principle of mass action is followed, 
but it only considers the antigen to stimulate immunity but ignores the immune damage 
caused by it. In  \cite{SN}, bilinear rate $\beta x(t)v(t)$ was replaced by infection 
rate $\frac{\beta x(t)v(t)}{1+\alpha v(t)}$ with saturation effect. Combined with
the actual situation, $py(t)z(t)$ was proposed in \cite{HZ,QJ, RW} to represent immune 
damage and included in the study of HIV dynamics. 

A large number of studies have shown that host immunity can be inhibited or even 
destroyed by certain pathogens, especially under the condition of excessive pathogen 
load, so antigens can impair immunity
(see,  \cite{RW,WS}). Combined with the situation of saturation incidence and 
intracellular delay, Xu \cite{X}  emphasized that the nonlinear function can better
reproduce the saturation reaction in the process of cellular immunity. 
Therefore, motivated by the ideas \cite{HZ, LX, WL, WS},
Deng and Xu \cite{DX} used the nonlinear function $\frac{py(t)z(t)}{1+\omega z(t)}$ 
instead of the bilinear function $py(t)z(t)$, and considered the model
\begin{equation}\label{1.3}
\begin{gathered}
\dot{x}(t)=\lambda-dx(t)-{\frac{\beta x(t)v(t)}{1+\alpha v(t)}},\\
\dot{y}(t)={\frac{\beta e^{-m\tau}x(t-\tau)v(t-\tau)}{1+\alpha v(t-\tau)}}
-ay(t)-{\frac{py(t)z(t)}{1+\omega z(t)}},\\
\dot{v}(t)=ky(t)- \mu v(t),\\
\dot{z}(t)=cy(t)-bz(t)-ny(t)z(t).
\end{gathered}
\end{equation}
Here, $\alpha, p, c, b, n$ are positive constants. $\alpha$ represents the constant 
rate of inhibition effect produced by the crowding effect of free virions;
$\tau$ is the time it takes from the entry of the virus into the susceptible cell to 
the production of new virions, and $e^{-m\tau}$ denotes the survival
probability of the infected cell from time $t-\tau$ to time $t$, $n$ is the rate of 
immune damage, $p$ represents the clearance rate of T cells to infected cells,
$c$ and $b$ represent the T cell production rate and death rate, respectively. 
$\omega$ represents the inhibition rate of CTL
immune response, and the meanings of other variables and parameters are the same as 
those in models \eqref{1.1} and \eqref{1.2}.

In \cite{DX}, the authors assumed that the production rate of susceptible cells in model 
\eqref{1.3} is a constant, that is, they did not take into account that susceptible 
cells  would also grow in the body. However, the results of the biological experiment 
indicate that the proliferation rate of CD4$^+$T cell is negatively
correlates with its absolute concentration (see, e.g.,\cite{BR}). 
Thus, many  scholars ( see, e.g., \cite {LGM, LS, PC, WLL, ZZZ})  recently 
proposed and studied that the  CD4$^+$T  proliferation is subject to logistic growth 
rate $rx(t)\big(1-\frac{x(t)}{x_{\rm max}}\big)$,  where $r$ is the rate of growth and
$x_{\rm max}$ is the maximum capacity of CD4$^+$ T in the human body.

From the biological perspective, we know that a viral infection or immune 
response is not transient in vivo, there may be time-lapse, and
a time delay is essential to explain a range of processes. 
It takes some time for infected cells to become active and produce virus particles; 
new virions also mature to become  infectious; and a time interval for antigen 
stimulation of immune cells. Therefore, it is necessary and important to consider time 
delays when modeling viral spread and  immune response.

 Based on the above analysis, in this paper, we consider the  
  delayed HIV infection model.
\begin{equation}\label{1.4}
\begin{gathered}
\dot{x}(t)=\lambda-dx(t)+rx(t)\Big(1-{\frac{x(t)}{x_{\rm max}}}\Big)-{\frac{\beta x(t)v(t)}{1+\alpha v(t)}},\\
\dot{y}(t)={\frac{\beta e^{-m\tau_1}x(t-\tau_1)v(t-\tau_1)}{1+\alpha v(t-\tau_1)}}-ay(t)-{\frac{py(t)z(t)}{1+\omega z(t)}},\\
\dot{v}(t)=k e^{-m\tau_2}y(t-\tau_2)-\mu v(t),\\
\dot{z}(t)=ce^{-m\tau_3}y(t-\tau_3)-bz(t)-ny(t)z(t).
\end{gathered}
\end{equation}
Here, uninfected target cells are assumed to be produced at a constant rate $\lambda$ 
and die at a per capita rate $d$. The proliferation of uninfected target cells is 
described by the logistic term $rx(t)\big(1-\frac{x(t)}{x_{\rm max}}\big)$, where
$r$ is the intrinsic mitosis rate and $x_{\rm max}$ is the carrying capacity. 
The delays during the processes of  viral production, and CTLs recruitment are $\tau_2$ 
and $\tau_3$, respectively. The term $ e^{-m\tau_2}$ describes the survival probability 
that start budding from activated infected cells at time $t$ and become free mature 
viruses at $\tau_2$ time later. The term $ e^{-m\tau_3}$ represents the survival rate 
of virus-specific CTLs during the delay
between cell encounters and subsequent recruitment. 
All parameters are positive constants. The other variables and parameters have the 
same meaning as model \eqref{1.3}.

The paper is structured as follows. 
In section 2, we study the well-posedness of model \eqref{1.4}. Then we give the 
basic reproduction number $\mathcal{R}_0$ and
the existence of the equilibria of  model \eqref{1.4}. 
In section 3, we show the basic reproduction number $\mathcal{R}_0$ determine
stability of the equilibria of model \eqref{1.4}. Then the locally  and global
asymptotic stability of the infection-free and the infection equilibrium
are derived.
In section 4, we illustrate the theoretical results with numerical simulations. 
Finally, in section 5, we provide a brief summary.


\section{Preliminaries}

In this section, we present some preliminary results including 
the well-posedness,  and the formula of
the basic reproduction number and existence of equilibria.


\subsection{Well-posedness of model \eqref{1.4}}

Let $C =C([-\tau,0];\mathbb{R}_{+}^{4})$ (where $\tau=\max \{\tau_1,\tau_2,\tau_3\}$) 
be the Banach space of continuous functions from $[-\tau, 0]$
to $\mathbb{R}$ with the norm $||\phi||=\max_{\theta\in [-\tau,0]}|\phi(\theta)|$ for 
any $\phi\in C$.
 For any given continuous function $u = (x, y, v, z): [-\tau,\sigma)$ with $\sigma>0$, 
we define
$$
u_t :=(x(t + \cdot), y(t + \cdot), v(t + \cdot), z(t))
\in C\times C\times C\times \mathbb{R}_+\quad \text{for  } t\in  [-\tau,\sigma).
$$
Let
$$
\mathbb{X}^+=C([-\tau,0], \mathbb{R}_+)\times C([-\tau,0], 
\mathbb{R}_+)\times C([-\tau,0], \mathbb{R}_+)\times \mathbb{R}_+.
$$
We first investigate the well-posedness of model \eqref{1.4}.

\begin{theorem}\label{T2.1}
For any $\phi \in \mathbb{X}^+ $, model (1.4) has a unique nonnegative solution 
$u(t,\phi)$  on $[0,\infty)$ with $u_0=\phi$, and $u_t(\phi)\in \mathbb{X}^+ $ for all
$t>0$.  Further, the solution semiflow $u_t: \mathbb{X}^+ \to \mathbb{X}^+$, $t>0$, 
has a compact global attractor.
\end{theorem}

\begin{proof}
For any $\phi=(\phi_1,\phi_2,\phi_3,\phi_4) \in \mathbb{X}^+$, we define
$$
f(\phi):=(f_1, f_2, f_3, f_4)(\phi),
$$
with
\begin{align*}
f_1(\phi)&=  \lambda-d\phi_1(0)+r\phi_1(0)\Big(1-\frac{\phi_1(0)}{x_{\rm max}}\Big)- \frac{{\beta }\phi_1(0)\phi_3(0)}{1+\alpha\phi_3(0)},\\
f_2(\phi)&=  \frac{\beta e^{-m\tau_1}\phi_1(-\tau_{1})\phi_3(-\tau_{1})}{1+\alpha\phi_3(-\tau_{1})}-a\phi_2(0)-\frac{p\phi_2(0)\phi_4(0)}{1+\omega\phi_4(0)},\\
f_3(\phi)&=ke^{-m\tau_2}\phi_2(-\tau_{2})-\mu\phi_3(0),\\
f_4(\phi)&=ce^{-m\tau_3}\phi_2(-\tau_3)-b\phi_4(0)-n\phi_2(0)\phi_4(0).
\end{align*}
Noting that $f(\phi)$ is continuous in  $\phi\in \mathbb{X}^+$ and
$f(\phi)$ is Lipschitz in $\phi$ on each compact subset of $\mathbb{X}^+$,  
then it implies that model \eqref{1.4} admits a unique solution $u(t,\phi)$ 
through $(0,\phi)$
on its maximal interval $[0, \sigma_{\phi})$  
(see \cite[Theorem 2.2.1 and Theorem 2.2.3]{HV}). Further, let
$\phi:= (\phi_1, \phi_2, \phi_3, \phi_4)\in \mathbb{X}^+$ be given, we can show 
that if $\phi_i(0)=0$ for $1\leq i\leq 4$, then $f_i(\phi)\geq 0$.
By \cite[Theorem 5.2.1 Remark 5.2.1]{S}, it follows that for any 
$(t, \phi)\in \mathbb{R}_+\times \mathbb{X}^+$, the unique solution $u(t,\phi)$
of  model \eqref{1.4} with $u_0=\phi$  satisfies $u_t(\phi)\in \mathbb{X}^+$ for all 
$t\in [0, \sigma_{\phi})$.

Next, we study the boundedness of solutions of model \eqref{1.4}. Let
 $$
 L(t)=x(t)+e^{m\tau_1}y(t+\tau_1).
 $$
Then,  in view of the first and second equations of \eqref{1.4}, we obtain
\begin{equation} \label{2.1}
 \begin{aligned}
 L'(t)&=\lambda-dx(t)+rx(t)\Big(1-\frac{x(t)}{x_{\rm max}}\Big)
  -\frac{\beta x(t)v(t)}{1+\alpha v(t)}   \\
 &\quad +\frac{\beta x(t)v(t)}{1+\alpha v(t)}-ae^{m\tau_1}y(t+\tau_1)-\frac{pe^{m\tau_1}y(t+\tau_1)z(t+\tau_1)}{1+\omega z(t+\tau_1)}  \\
 &\leqslant\lambda-dx(t)+rx(t)\Big(1-\frac{x(t)}{x_{\rm max}}\Big)-ae^{m\tau_1}y(t+\tau_1)   \\
% &=-dx(t)-ae^{m\tau_1}y(t+\tau_1)-\frac{r}{x_{\rm max}}(x(t)-\frac{x_{\rm max}}{2})^{2}+\lambda+{\frac{x_{\rm max}}{4}}r\\
 &\leqslant\lambda+\frac{1}{4}rx_{\rm max}-dx(t)-ae^{m\tau_1}y(t+\tau_1)   \\
 &\leqslant \lambda+\frac{1}{4}rx_{\rm max}-\delta L(t),
 \end{aligned}
\end{equation}
where $\delta=\min\{d,a\}$. By the comparison theorem,  we see $L(t)$ is bounded
on $[0, \sigma_{\phi})$, it implies $x(t)$ and $y(t)$ are all bounded on
$[0, \sigma_{\phi})$.  As a result, we assume that there is $M>0$ such that
$|y(t)|\leq M$ for all $t\in [0, \sigma_{\phi})$. It then follows from  the third and
the last equation of  \eqref{1.4} that we have
 \begin{gather}
 \frac{dv(t)}{dt}\leqslant k Me^{-m\tau_2}-\mu v(t),  \label{2.2} \\
 \frac{dz(t)}{dt}\leqslant cMe^{-m\tau_3}-bz(t).\label{2.3}
 \end{gather}
Hence, both $v(t)$ and $z(t)$ are bounded on $[0, \sigma_{\phi})$, and then
\cite[Theorem 2.3.1]{HV} follows that $\sigma=\infty$.
By the differential inequalities \eqref{2.1}-\eqref{2.3} and the comparison theorem,
it then follows that solutions of model \eqref{1.4} are uniformly bounded and
ultimately bounded in $\mathbb{X}^+$. Consequently,
\cite [Corollary 3.6.2 and Theorem 4.5.2]{HV}
imply the existence of a compact global attractor for the solution semiflow of model
\eqref{1.4} on $\mathbb{X}^+$.
\end{proof}


\subsection{Basic reproductive number and equilibria}

In this subsection, we  give the basic reproduction number for model \eqref{1.4}, 
and then establish the existence of equilibria of model \eqref{1.4}.  
Obviously,   model \eqref{1.4} always has an infection-free equilibrium
 $E_0=(x_0,0,0,0)$, where
\begin{equation} \label{2.4}
 x_0=\frac{x_{\rm max}}{2r} \Big(r-d+\sqrt{(r-d)^{2}
 +\frac{4r\lambda}{x_{\rm max}}}\, \Big).
\end{equation}
Linearizing \eqref{1.4} at $E_0$, we obtain the following three infection related
equations for variables $y$, $v$ and $z$ satisfied
 \begin{align*}
\dot{y}(t)&= \beta e^{-m\tau_1}x_0v(t-\tau_1)-ay(t),\\
\dot{v}(t)&= k e^{-m\tau_2}y(t-\tau_2)-\mu v(t),\\
\dot{z}(t)&= ce^{-m\tau_3}y(t-\tau_3)-bz(t).
\end{align*}
Let $u_1$, $u_2$ and $u_3$ be the number of infected cells, HIV and CTL cell at $t=0$,
respectively. Then the remaining numbers of infected cells
 HIV and CTL  cell at time $t$ are given by
$$
u_1(t)=u_1e^{-at},\quad u_2(t)=u_2e^{-\mu t}, \quad u_3(t)=u_3e^{-bt}.
$$
The total numbers of newly infected cells, produced HIV and CTL  cell are
\begin{align*}
\bar{u}_1&=  \int_{\tau_{1}}^{\infty}\beta e^{-m\tau_1}x_0u_2(t-\tau_1)dt=\frac{\beta }{\mu}e^{-m\tau_1}x_0u_2,\\
\bar{u}_2&=  \int_{\tau_{2}}^{\infty}k e^{-m\tau_2} u_1(t-\tau_2)dt=\frac{k}{a}e^{-m\tau_2}u_1,\\
\bar{u}_3&= \int_{\tau_{3}}^{\infty}ce^{-m\tau_3} u_1(t-\tau_3)dt=\frac{c}{b}e^{-m\tau_3}u_1,
\end{align*}
which can be rewritten as
\begin{align*}
\begin{array}{ll}
\begin{pmatrix} \bar{u}_1\\
\bar{u}_2\\
\bar{u}_3\\
\end{pmatrix}=\mathcal{M}_0
\begin{pmatrix}
 {u}_1\\
 {u}_2\\
 {u}_3\\
 \end{pmatrix},
 \quad \text{where } 
 \mathcal{M}_0=\begin{pmatrix}
 0&\frac{\beta }{\mu} e^{-m \tau_1}x_0&0\\
 \frac{k}{a} e^{-m\tau_2}&0&0\\
 \frac{c}{b}e^{-m\tau_3}&0&0\\
 \end{pmatrix}.
 \end{array}
 \end{align*}
Then, we can see the matrix $\mathcal{M}_0$ is the next infection operator \cite{DW}.
 As usual, the spectral radius of $\mathcal{M}_0$ is called the basic reproduction
number $\mathcal{R}_0$,  which is
\[
 \mathcal{R}_0=\sqrt{{\frac{k\beta x_0}{a\mu}e^{-m(\tau_1+\tau_2)}}}.
\]

Next, we shall establish the existence of  the infection equilibrium of model 
\eqref{1.4}. The infection equilibrium denoted by $E^{*}=(x^{*},y^{*},v^{*},z^{*})$ 
satisfies the system
 \begin{equation} \label{2.5}
 \begin{gathered}
 h(x)-{\frac{\beta xv}{1+\alpha v}}=0,\\
 {\frac{\beta e^{-m\tau_1}xv}{1+\alpha v}}-ay-{\frac{pyz}{1+\omega z}}=0,\\
 k e^{-m\tau_2}y-\mu v=0,\\
 ce^{-m\tau_3}y-bz-nyz=0,
\end{gathered}
\end{equation}
where $h(x)= \lambda-dx(t)+rx(t)(1-\frac{x(t)}{x_{\rm max}})$.
It follows from  the first equation of \eqref{2.5} that
$$
\frac{r}{x_{\rm max}}x^2+\Big(d-r+\frac{\beta v}{1+\alpha v}\Big)x -\lambda=0.
$$
Clearly, the positive root $x^*$ of the above equation is
$$
x^*=\frac{x_{\rm max}}{2r}
\Big(r-d-\frac{\beta v}{1+\alpha v}+\sqrt{\Big(r-d-\frac{\beta v}{1+\alpha v}
\Big)^{2}+\frac{4r\lambda}{x_{\rm max}}}\;\Big)<x_0.
$$
Obviously, $\beta x^*-\alpha h(x^*)\not=0$.  From the first equation
of \eqref{2.5} again, we obtain
\[
v=\frac{h(x)}{\beta x-\alpha h(x)}.
\]
By the third and fourth equations of  \eqref{2.5}, respectively,  we obtain
$$
y=\frac{\mu }{k}e^{m\tau_2}v, \;\;z=ce^{-m \tau_3}\frac{y}{b+ny}.
$$
Substituting $y=\frac{\mu}{k}e^{m\tau_2}v$ into $z=ce^{-m \tau_3}\frac{y}{b+ny}$,
we have
\[
z=\frac{Bv}{1+Cv}, \quad \text{where}\quad B=\frac{c\mu}{kb}
e^{-m\tau_3+m\tau_2},\;\; C=\frac{n \mu}{kb} e^{m\tau_2}.
\]
Putting $z=\frac{Bv}{1+Cv}$ into the second equation of \eqref{2.5}, we obtain
\[
h(x)=\Big(a+\frac{pBv}{1+Cv+ \omega Bv}\Big)Av, \quad \text{where}\quad
A=\frac{\mu}{k}e^{m(\tau_1+\tau_2)},
\]
that is,
\[
h(x)(1+Cv+\omega Bv)=A(a+(aC+a\omega B+pB)v)v.
\]
Substituting $v=\frac{h(x)}{\beta x-\alpha h(x)}$ into the above equation
and noting that $h(x)\not=0$, we obtain
\[
\Big[\beta x-\alpha h(x)+(C+\omega B)h(x)\Big]
(\beta x-\alpha h(x))=A\Big[a(\beta x-\alpha h(x))+(aC+a\omega B+pB)h(x)\Big].
\]
We define
\begin{equation} \label{2.6}
\begin{aligned}
H(x)&=  \Big[\beta x-\alpha h(x)+(C+\omega B)h(x)\Big](\beta x-\alpha h(x))    \\
&\quad -A\Big[a(\beta x-\alpha h(x))+(aC+a\omega B+pB)h(x)\Big]\quad
 \text{for } 0<x<x_0.
\end{aligned}
\end{equation}

Now, we claim that there is a positive root $x^*$ in the following equation
\[
H(x)=0, \quad  0<x<x_0.
\]
In fact, it is easy to see the equation  $\alpha h(x)-\beta x=0$  has the positive root
$$
\tilde{x}=\frac{x_{\rm max}}{2r}\Big(r-d-\frac{\beta }{\alpha}
+\sqrt{\Big(r-d-\frac{\beta}{\alpha}\Big)^{2}+\frac{4r\lambda}{x_{\rm max}}}\Big),
$$
and  $0<\tilde{x}<x_0$.

Since $h(\tilde{x}_0)-\beta \tilde{x}_0=0$, we have
\[
H(\tilde{x}_0)=-A\beta (aC+a\omega B+pB)\tilde{x}_0<0.
\]
And also, if $\mathcal{R}_0>1$, noting that $h(x_0)=0$, we then obatin
\[
H(x_0)=(\beta x_0)^2-Aa\beta x_0=\frac{a\mu}{k}
\beta x_0e^{m(\tau_1+\tau_2)}(\mathcal{R}_0^{2}-1)>0.
\]
Therefore, there is a $x^*\in (\tilde{x}_0, x_0)$ such that $H(x^*)=0$.

Summarizing the above discussion, we  establish the existence of equilibria of 
model \eqref{1.4}.

\begin{theorem}\label{T2.2}
(1) Model \eqref{1.4} always has the infection-free equilibrium $E_0(x_0,0,0,0)$ with 
$x_0$ defined by \eqref{2.1}.

(2) Model \eqref{1.4} has the infection equilibrium $E^{*}=(x^{*}, y^{*}, v^{*}, z^{*})$ 
if $\mathcal{R}_0>1$, where  $x^*$ is the positive root of $H(x)=0$, and
\[
v^*=\frac{h(x^*)}{\beta x^*-\alpha h(x^*)},\quad 
y^*=\frac{u}{k}e^{m\tau_2}v^*,\quad z^*=\frac{u}{k}e^{m\tau_2}v^*,
\]
where  $H(x)$ is defined by \eqref{2.6}.
\end{theorem}

\section{Stability analysis}

Here  we  establish the local and global stabilities of the equilibria of 
model \eqref{1.4}.

\subsection{Local stability of equilibria of model \eqref{1.4}}

By analyzing the distribution of the roots of the corresponding characteristic equation, 
we  study the local stability of the infection-free and infection
equilibrium of model \eqref{1.4}. First, we have the following result.

\begin{theorem}\label{T3.1}
Let $\tau_{i}\geqslant0,\ i=1,2,3$. The infection-free equilibrium  $E_0$ is 
locally asymptotically stable for $\mathcal{R}_0<1$ and it is unstable for 
$\mathcal{R}_0>1$.
\end{theorem}

\begin{proof}
Linearizing \eqref{1.4} at $E_0$, we can see the Jacobian matrix at $E_0$ as follows
\[
J_{E_0}=\begin{pmatrix}
-d+r-\frac{2r}{x_{\rm max}}x_0&0&-\beta x_0&0\\
0&-a&\beta x_0e^{-(\xi+m)\tau_1}& 0\\
0&ke^{-(\xi+m)\tau_2}&-\mu&0\\ 0&ce^{-(\xi+m)\tau_3}&0&-b
\end{pmatrix}.
\]
Then the characteristic equation of $J_{E_0}$ is
\begin{equation}\label{3.1}
(\xi+b)\Big(\xi+d-r+{\frac{2r}{x_{\rm max}}x_0}\Big)\Big((\xi+a)(\xi+\mu)
-k\beta x_0e^{-(\xi+m)(\tau_1+\tau_2)}\Big)=0.
\end{equation}
It is clear that equation \eqref{3.1} has two negative roots
\[
\xi_1=-b \quad \text{and}\quad  \xi_2=-\Big(d-r+\frac{2r}{x_{\rm max}}x_0\Big),
\]
since
$$
d-r+{\frac{2r}{x_{\rm max}}x_0}=\sqrt{(r-d)^{2}+\frac{4r\lambda}{x_{\rm max}}}>0.
$$
Hence, the remaining roots of equation \eqref{3.1} are determined by the 
equation
$$
\kappa(\xi)=:(\xi+a)(\xi+\mu)-k\beta x_0e^{-(\xi+m)(\tau_1+\tau_2)}= 0.
$$	
Obviously, the equation $\kappa(\xi)=0$ is equivalent to the equation
\begin{equation}\label{3.2}
1=\frac{k\beta x_0 }{(\xi+a)(\xi+\mu)}e^{-(\xi+m)(\tau_1+\tau_2)}.
\end{equation}	
Let $\xi=x+iy $ with $x>0$ is a solution of equation \eqref{3.2}, and the modulus 
on both sides of equation \eqref{3.2}, we have
\begin{align*}
1=\Big|\frac{k\beta x_0}{(\xi+a)(\xi+\mu)}e^{-(\xi+m)(\tau_1+\tau_2)}\Big|
\leqslant \Big|\frac{k\beta x_0 }{a\mu}\Big||e^{-m(\tau_1+\tau_2)}|= \mathcal{R}_0^{2}<1,
\end{align*}
which implies $1\leq \mathcal{R}_0<1$, a contradiction. 
Hence all roots of equation \eqref{3.2} have
negative real parts if $\mathcal{R}_0<1$, it follows that $E_0$ is locally 
asymptotically stable for $\mathcal{R}_0<1$.
	
Noting that
$$
\kappa(0)=a\mu(1-\mathcal{R}_0^{2})<0, \quad 
\lim_{\xi\to+\infty}\kappa(\xi)=+\infty.
$$
Hence, the equation $\kappa(\xi)=0$ has a positive root if $\mathcal{R}_0>1$. 
It follows that $E_0$ is unstable if $\mathcal{R}_0>1$.	
\end{proof}

\begin{theorem}\label{T3.2}
Let $\tau_{i}\geqslant0$, $i=1,2$, and $\tau_3=0$. The infection equilibrium $E^{*}$ 
is locally asymptotically stable for $\mathcal{R}_0>1$.
\end{theorem}

\begin{proof}
Linearizing the system \eqref{1.4} at $E^*$, we can get the Jacobian matrix at $E^{*}$  
given by	
\[
J_{E^{*}}
=\begin{pmatrix}
-d+r-\frac{2r}{x_{\rm max}}x^*-\frac{\beta v^{*}}{1+\alpha v^{*}}&0&-\frac{\beta x^{*}}{(1+\alpha v^{*})^{2}}&0\\[4mm]
\frac{\beta v^{*}}{1+\alpha v^{*}}e^{-(\xi+m)\tau_1}&-a-\frac{pz^{*}}{1+\omega z^{*}}&\frac{\beta x^{*}}{(1+\alpha v^{*})^{2}}e^{-(\xi+m)\tau_1}
&-\frac{pz^{*}}{(1+\omega z^{*})^{2}}\\[4mm]
 0&ke^{-(\xi+m)\tau_2}&-\mu&0\\[4mm]
 0&c-nz^{*}&0&-b-ny^{*}
 \end{pmatrix}.
\]
Thus, we obtain the following characteristic equation of \eqref{1.4} at $E^{*}$ 
as follows
\begin{equation} \label{3.3}
\begin{aligned}
&\Big(\xi+d-r+{\frac{2r}{x_{\rm max}}x^{*}}+\frac{\beta v^{*}}{1+\alpha v^{*}}\Big)(\xi+\mu)  \\[5mm]
&\times\Big(\Big(\xi+a+{\frac{pz^{*}}{1+\omega z^{*}}}\Big)
(\xi+b+ny^{*})+{\frac{py^{*}}{(1+\omega z^{*})^{2}}}(c-nz^{*})\Big)    \\[5mm]
&=\frac{k\beta x^*}{(1+\alpha v^*)^2}\Big(\xi+d-r+\frac{2r}{x_{\rm max}}x^{*}\Big)(\xi+b+ny^{*})e^{-(\xi+m)(\tau_1+\tau_2)}.
\end{aligned}
\end{equation}
Next, we claim all roots of equation \eqref{3.3} have negative real parts when
$\mathcal{R}_0>1$. Otherwise, equation \eqref{3.3} has a root $\xi=x+iy $ with
$x>0$, and note that
\begin{equation*}
c-nz^{*}=\frac{bz^{*}}{y^{*}}, \quad y^{*}=\frac{u}{k}e^{m\tau_2}v^{*},\quad
\frac{\beta x^{*}v^*}{(1+\alpha v^{*})y^{*} }=a+\frac{pz^{*}}{1+\omega z^{*}}.
\end{equation*}
Thus,
\begin{equation} \label{3.4}
\Big|\xi+d-r+\frac{2r}{x_{\rm max}}x^{*}+\frac{\beta v^{*}}{1+\alpha v^{*}}\Big|
>\Big|\Big(\xi+d-r+\frac{2r}{x_{\rm max}}x^{*}\Big)e^{-\xi(\tau_1+\tau_2)}\Big|,
\end{equation}
and
\begin{equation} \label{3.5}
\begin{aligned}
&\Big|(\xi+\mu)\Big(\Big(\xi+a+{\frac{pz^{*}}{1+\omega z^{*}}}\Big)
 (\xi+b+ny^{*})+{\frac{py^{*}}{(1+\omega z^{*})^{2}}}(c-nz^{*})\Big)\Big|  \\
&= \Big|(\xi+\mu)\Big(\Big(\xi+a+{\frac{pz^{*}}{1+\omega z^{*}}}\Big)
 (\xi+b+ny^{*})+{\frac{bpz^{*}}{(1+\omega z^{*})^{2}}}\Big)\Big|  \\
&\geqslant \Big|\mu(\xi+b+ny^{*})\Big(a+\frac{pz^*}{1+\omega z^*}\Big)\Big|    \\
&\geqslant \Big|(\xi+b+ny^{*})\frac{k\beta x^{*}}{1+\alpha v^{*}}
 e^{-m(\tau_1+\tau_2)}\Big|     \\
&> \Big|(\xi+b+ny^{*})\frac{k\beta x^{*}}{(1+\alpha v^{*})^{2}}
 e^{-m(\tau_1+\tau_2)}\Big|.
\end{aligned}
\end{equation}
In view of inequalities \eqref{3.4} and \eqref{3.5}, we obtain that $\xi=x+iy$ with $x>0$
is not the root of \eqref{3.3}, which leads that
equation \eqref{3.3} can not have any root with positive real part.
Hence, the  infection equilibrium $E^{*}=(x^{*},y^{*},v^{*},z^{*})$ of model \eqref{1.4}
is locally asymptotically stable for $\mathcal{R}_0>1$.
\end{proof}


\subsection{Global stability of equilibria of model \eqref{1.4}}

Constructing Lyapunov functionals, we study the global asymptotic stability of 
the  equilibria of model \eqref{1.4}.
To this end,  we introduce the fundamental function
$$
g(x)=x-1-\ln x \quad \text{for }  x>0.
$$
Clearly, the function $g(x) \geqslant$0 for all $x>0$, and $g(x)=0$ if and only 
if $x=1$.

We start with the globally  asymptotic stability of the infection-free  
equilibrium $E_0$.

\begin{theorem}\label{T3.3}
Let $\tau_{i}\geqslant0,\ i=1,2,3$.  If $\mathcal{R}_0\leq 1$, then the 
infection-free equilibrium  $E_0$ of model \eqref{1.4} is globally asymptotically stable.
\end{theorem}

\begin{proof}
We define the Lyapunov functional 
$$
M(t)=M_{1}(t)+M_{2}(t),
$$
where
\begin{gather*}
M_{1}(t)=x_0 g\Big(\frac{x(t)}{x_0}\Big)
+e^{m\tau_1}y(t)+\frac{\beta x_0}{\mu}e^{m(\tau_1+\tau_2)}v(t)
+e^{m(\tau_1+\tau_3)}(1-\mathcal{R}_0^{2})z(t),
\\
M_{2}(t)=\int _{t-\tau_{1}}^{t}\frac{\beta x(s)v(s)}{1+\alpha v(s)}ds
+\mathcal{R}_0^{2}ae^{-m\tau_2}\int_{t-\tau_{2}}^{t}y(s)ds
+ae^{m\tau_1}(1-\mathcal{R}_0^{2}) \int _{t-\tau_{3}}^{t}y(s)ds.
\end{gather*}
It can be easily verified that $M(t)\geq 0$ with  $M(t)= 0$ if and only if 
$x=x_0$, $y=0$, $v=0$ and $z=0$.
 Differentiating $M(t)$ along  any positive solution of model \eqref{1.4}, we have
\begin{align*}
M_{1}'(t)
&= -\frac{(x_0-x(t))^{2}}{x(t)}\Big(d-r+\frac{rx_0}{x_{\rm max}}
 +\frac{rx(t)}{x_{\rm max}}\Big)-\frac{\beta x(t)v(t)}{1+\alpha v(t)}
 +\frac{\beta x(t)v(t)}{1+\alpha v(t)}\frac{x_0}{x(t)} \\
&\quad +\frac{\beta x(t-\tau_{1})v(t-\tau_{1})}{1+\alpha v(t-\tau_{1})}
 -ae^{m\tau_1}y(t)-\frac{pe^{m\tau_{1}} y(t)z(t)}{1+\omega z(t)}\\
&\quad +\frac{k\beta x_0}{\mu}e^{m\tau_1}y(t-\tau_{2})-\beta x_0e^{m(\tau_1+\tau_2)}v(t) \\
&\quad +ae^{m\tau_1}(1-\mathcal{R}_0^{2})y(t-\tau_{3})
 -\frac{a}{c}e^{m(\tau_1+\tau_3)}(1-\mathcal{R}_0^{2})(bz(t)+ny(t)z(t)),
\end{align*}
and
\begin{align*}
M_{2}'(t)
&= \frac{\beta x(t)v(t)}{1+\alpha v(t)}
 -\frac{\beta x(t-\tau_{1})v(t-\tau_{1})}{1+\alpha v(t-\tau_{1})}
+\mathcal{R}_0^{2}ae^{-m\tau_2}y(t)-\mathcal{R}_0^{2}ae^{-m\tau_2}y(t-\tau_{2})\\
&\quad +ae^{m\tau_1}(1-\mathcal{R}_0^{2})y(t)-ae^{m\tau_1}
 (1-\mathcal{R}_0^{2})y(t-\tau_{3}).
\end{align*}
Thus,  we obtain
\begin{align*}
M'(t)&= -\frac{(x_0-x(t))^{2}}{x(t)}\Big(d-r+\frac{rx_0}{x_{\rm max}}
 +\frac{rx(t)}{x_{\rm max}}\Big)+\frac{\beta x(t)v(t)}{1+\alpha v(t)}\frac{x_0}{x(t)}
-\frac{pe^{m\tau_{1}} y(t)z(t)}{1+\omega z(t)}\\
&\quad -\frac{a}{c}e^{m(\tau_1+\tau_3)}(1-\mathcal{R}_0^{2})(bz(t)+ny(t)z(t))
 +\mathcal{R}_0^{2}ae^{-m\tau_2}y(t)-\mathcal{R}_0^{2}ae^{-m\tau_1}y(t)\\
&= -\frac{(x_0-x(t))^{2}}{x(t)}\Big(d-r+\frac{rx_0}{x_{\rm max}}
 +\frac{rx(t)}{x_{\rm max}}\Big)-\beta x_0v(t)\Big(e^{m(\tau_1+\tau_2)}
 -\frac{1}{1+\alpha v(t)}\Big)\\
&\quad -\mathcal{R}_0^{2}ae^{-m\tau_1}y(t)(1-e^{-m(\tau_1+\tau_2)})
 -\frac{pe^{m\tau_{1}} y(t)z(t)}{1+\omega z(t)}\\
&\quad -\frac{a}{c}e^{m(\tau_1+\tau_3)}(1-\mathcal{R}_0^{2})(bz(t)+ny(t)z(t)).
\end{align*}
Note that $d-r+\frac{rx_0}{x_{\rm max}}>0$. 
If $\mathcal{R}_0<1$, then $M'(t)\leqslant0$ and the equality holds if and only if 
 $x(t)=x_0$, $y(t)=0$, $v(t)=0$ and $z(t)=0$. 
Hence, the largest invariant set in the set $\{(x,y,v,z)\!: M'(t)=0\}$
is the singleton set $E_0$. Therefore, by the LaSalle's invariance 
principle \cite{HV}, the infection-free equilibrium $E_0$ is
globally asymptotically stable when $\mathcal{R}_0<1$.

Now, if $\mathcal{R}_0=1$, we obtain
\begin{align*}
M'(t)
&= -\frac{(x_0-x(t))^{2}}{x(t)}\Big(d-r+\frac{rx_0}{x_{\rm max}}
 +\frac{rx(t)}{x_{\rm max}}\Big)-\beta x_0v(t)\Big(e^{m(\tau_1+\tau_2)}
 -\frac{1}{1+\alpha v(t)}\Big)\\
&\quad -ae^{-m\tau_1}y(t)(1-e^{-m(\tau_1+\tau_2)})
 -\frac{pe^{m\tau_{1}} y(t)z(t)}{1+\omega z(t)},
\end{align*}
which follows $M'(t)\leqslant0$ and the equality holds if and only if  $x(t)=x_0$, 
$y(t)=0$, $v(t)=0$. Then one can easily show that the largest invariant set of
$\{(x,y,v,z): M'(t)=0\}$ is the singleton $\{E_0\}$. Hence, using LaSalle's 
invariance principle \cite{HV}, we also get $E_0$ is globally asymptotically
stable when $\mathcal{R}_0=1$.
\end{proof}


\begin{theorem}\label{T3.4}
Let $\tau_{i}\geqslant0$, $i=1,2$, and $\tau_3=0$.  If $\mathcal{R}_0> 1$ and that 
$x_{\rm max}(r-d)<rx^{*}$, then the infection equilibrium $E^{*}$ of model 
\eqref{1.4} is globally asymptotically stable.
\end{theorem}

\begin{proof}
We define the Lyapunov function 
$$
N(t)=N_{1}(t)+N_{2}(t),
$$
where
\begin{gather*}
N_{1}(t)=x^{*} g\Big(\frac{x(t)}{x^{*}}\Big)
 +e^{m\tau_1}g\Big(\frac{y(t)}{y^{*}}\Big)+\frac{\beta x^{*}v^{*}}{(1+\alpha v^{*})ky^{*}}e^{m\tau_2}v^{*}
g\Big(\frac{v(t)}{v^{*}}\Big)+\frac{py^{*}z^{*}}{b(1+\omega z^{*})^{2}}
e^{m\tau_1}g\Big(\frac{z(t)}{z^{*}}\Big), \\
N_{2}(t)=\frac{\beta x^{*}v^{*}}{1+\alpha v^{*}}\int _{t-\tau_{1}}^{t}
 g\Big(\frac{x(s)v(s)(1+\alpha v^{*})}{x^{*}v^{*}(1+\alpha v(s))}\Big)ds
+\frac{\beta x^{*}v^{*}}{1+\alpha v^{*}}\int _{t-\tau_{2}}^{t}
 g\Big(\frac{y(s)}{y^{*}}\Big)ds.
\end{gather*}
Then, differentiating $N(t)$ along any positive solution of model \eqref{1.4} 
for $t\geq 0$, we obtain
\begin{align*}
N_{1}'(t)
&= \Big(1-\frac{x^{*}}{x(t)}\Big)\Big(\lambda-dx(t)
 -{\frac{\beta x(t)v(t)}{1+\alpha v(t)}}\Big)\\
&\quad +e^{m\tau_1}\Big(1-\frac{y^{*}}{y(t)}\Big)
 \Big({\frac{\beta e^{-m\tau}x(t-\tau)v(t-\tau)}{1+\alpha v(t-\tau)}}-ay(t)
 -{\frac{py(t)z(t)}{1+\omega z(t)}}\Big)\\
&\quad +\frac{\beta x^{*}v^{*}}{1+\alpha v^{*}}\frac{1}{ky^{*}}e^{m\tau_2}
 \Big(1-\frac{v^{*}}{v(t)}\Big)(ke^{-m\tau_2}y(t-\tau_2)-\mu v(t))\\
&\quad +\frac{py^{*}}{b(1+\omega z^{*})^{2}}e^{m\tau_1}
 \Big(1-\frac{x^{*}}{x(t)}\Big)(cy(t)-bz(t)-ny(t)z(t))\\
&= \Big(1-\frac{x^{*}}{x(t)}\Big)
 \Big(\lambda-dx(t)-{\frac{\beta x(t)v(t)}{1+\alpha v(t)}}\Big)
 +\frac{\beta x(t-\tau_{1})v(t-\tau_{1})}{1+\alpha v(t-\tau_{1})}-ae^{m\tau_1}y(t) \\
&\quad -\frac{pe^{m\tau_{1}} y(t)z(t)}{1+\omega z(t)}
 -\frac{\beta x(t-\tau_{1})v(t-\tau_{1})}{1+\alpha v(t-\tau_{1})}\frac{y^{*}}{y(t)}
 +ae^{m\tau_1}y^{*}+\frac{pe^{m\tau_{1}} y^{*}z(t)}{1+\omega z(t)}\\
&\quad +\frac{\beta x^{*}v^{*}}{1+\alpha v^{*}}\frac{y(t-\tau_{2})}{y^{*}}
 -\frac{\beta x^{*}v^{*}}{1+\alpha v^{*}}\frac{\mu}{ky^{*}}e^{m\tau_2}v(t)
 -\frac{\beta x^{*}v^{*}}{1+\alpha v^{*}}\frac{y(t-\tau_{2})}{y^{*}}\frac{v^{*}}{v(t)}\\
&\quad +\frac{\beta x^{*}v^{*}}{1+\alpha v^{*}}\frac{\mu v^{*}}{ky^{*}}e^{m\tau_2}
 +\frac{py^{*}}{b(1+\omega z^{*})^{2}}e^{m\tau_1}\frac{z(t)-z^{*}}{z(t)}(cy(t)
 -bz(t)-ny(t)z(t)),
\end{align*}
and
\begin{align*}
N_{2}'(t)
&= \frac{\beta x(t)v(t)}{1+\alpha v(t)}
 -\frac{\beta x(t-\tau_{1})v(t-\tau_{1})}{1+\alpha v(t-\tau_{1})}
 +\frac{\beta x^{*}v^{*}}{1+\alpha v^{*}}
 \ln\frac{x(t-\tau_{1})v(t-\tau_{1})(1+\alpha v(t))}{x(t)v(t)(1+\alpha v(t-\tau_{1}))}\\
&\quad + \frac{\beta x^{*}v^{*}}{1+\alpha v^{*}}\frac{y(t)}{y^{*}}
-\frac{\beta x^{*}v^{*}}{1+\alpha v^{*}}\frac{y(t-\tau_{2})}{y^{*}}
+\frac{\beta x^{*}v^{*}}{1+\alpha v^{*}}\ln\frac{y(t-\tau_{2})}{y(t)}.
\end{align*}
Noting that the infection equilibrium $E^{*}(x^{*}, y^{*}, v^{*}, z^{*})$ 
satisfies the equations
\begin{gather*}
\lambda=dx^{*}-rx^{*}\Big(1-{\frac{x^{*}}{x_{\rm max}}}\Big)
+{\frac{\beta x^{*}v^{*}}{1+\alpha v^{*}}}, \\
{\frac{\beta e^{-m\tau_1}x^{*}v^{*}}{1+\alpha v^{*}}}=ay^{*}
 +{\frac{py^{*}z^{*}}{1+\omega z^{*}}},\quad
k e^{-m\tau_2}y^{*}=\mu v^{*},\quad
 cy^{*}=bz^{*}+ny^{*}z^{*},
\end{gather*}	
we then have
\begin{align*}
N'(t)
&= -\frac{(x^{*}-x(t))^{2}}{x(t)}\Big(d-r+\frac{rx^{*}}{x_{\rm max}}+\frac{rx(t)}{x_{\rm max}}\Big)
 +\frac{\beta x^{*}v^{*}}{1+\alpha v^{*}}-\frac{\beta x(t)v(t)}{1+\alpha v(t)}
 -\frac{\beta x^{*}v^{*}}{1+\alpha v^{*}}\frac{x^{*}}{x(t)} \\
&\quad +\frac{\beta x(t)v(t)}{1+\alpha v(t)}\frac{x^{*}}{x(t)}
 +\frac{\beta x(t-\tau_{1})v(t-\tau_{1})}{1+\alpha v(t-\tau_{1})}
 -\frac{\beta x^{*}v^{*}}{1+\alpha v^{*}}\frac{y(t)}{y^{*}}
 +\frac{py^{*}z^{*}}{1+\omega z^{*}}e^{m\tau_1}\frac{y(t)}{y^{*}} \\
&\quad -\frac{pe^{m\tau_{1}} y(t)z(t)}{1+\omega z(t)}
 -\frac{\beta x(t-\tau_{1})v(t-\tau_{1})}{1+\alpha v(t-\tau_{1})}\frac{y^{*}}{y(t)}
 +\frac{\beta x^{*}v^{*}}{1+\alpha v^{*}}-\frac{py^{*}z^{*}}{1+\omega z^{*}}e^{m\tau_1} \\
&\quad +\frac{pe^{m\tau_{1}} y^{*}z(t)}{1+\omega z(t)}
 +\frac{\beta x^{*}v^{*}}{1+\alpha v^{*}}\frac{y(t-\tau_{2})}{y^{*}}
 -\frac{\beta x^{*}v^{*}}{1+\alpha v^{*}}\frac{\mu}{ky^{*}}e^{m\tau_2}v(t)
 -\frac{\beta x^{*}v^{*}}{1+\alpha v^{*}}\frac{y(t-\tau_{2})}{y^{*}}\frac{v^{*}}{v(t)} \\
&\quad +\frac{\beta x^{*}v^{*}}{1+\alpha v^{*}}+\frac{pe^{m\tau_1}y^{*}(z(t)
 -z^{*})}{bz(t)(1+\omega z^{*})^{2}}\Big(ny(t)(z^{*}-z(t))+\frac{by(t)z^{*}
 -by^{*}z(t)}{y^{*}}\Big)\\
&\quad +\frac{\beta x(t)v(t)}{1+\alpha v(t)}
 -\frac{\beta x(t-\tau_{1})v(t-\tau_{1})}{1+\alpha v(t-\tau_{1})}
 +\frac{\beta x^{*}v^{*}}{1+\alpha v^{*}}\ln\frac{x(t-\tau_{1})
  v(t-\tau_{1})(1+\alpha v(t))}{x(t)v(t)(1+\alpha v(t-\tau_{1}))}\\
&\quad +\frac{\beta x^{*}v^{*}}{1+\alpha v^{*}}\frac{y(t)}{y^{*}} 
 -\frac{\beta x^{*}v^{*}}{1+\alpha v^{*}}\frac{y(t-\tau_{2})}{y^{*}}
 +\frac{\beta x^{*}v^{*}}{1+\alpha v^{*}}\ln\frac{y(t-\tau_{2})}{y(t)}\\
&= -\frac{(x^{*}-x(t))^{2}}{x(t)}\Big(d-r+\frac{rx^{*}}{x_{\rm max}}
 +\frac{rx(t)}{x_{\rm max}}\Big)+\frac{\beta x^{*}v^{*}}{1+\alpha v^{*}}
 -\frac{\beta x^{*}v^{*}}{1+\alpha v^{*}}\frac{x^{*}}{x(t)}
 +\frac{\beta x(t)v(t)}{1+\alpha v(t)}\frac{x^{*}}{x(t)} \\
&\quad -\frac{\beta x(t-\tau_{1})v(t-\tau_{1})}{1+\alpha v(t-\tau_{1})}
 \frac{y^{*}}{y(t)}+\frac{\beta x^{*}v^{*}}{1+\alpha v^{*}}
-\frac{\beta x^{*}v^{*}}{1+\alpha v^{*}}\frac{v(t)}{v^{*}}
 -\frac{\beta x^{*}v^{*}}{1+\alpha v^{*}}\frac{y(t-\tau_{2})}{y^{*}}\frac{v^{*}}{v(t)} \\
&\quad +\frac{\beta x^{*}v^{*}}{1+\alpha v^{*}}
 +\frac{py^{*}z^{*}}{1+\omega z^{*}}e^{m\tau_1}\frac{y(t)}{y^{*}}
 -\frac{pe^{m\tau_{1}} y(t)z(t)}{1+\omega z(t)}
 -\frac{py^{*}z^{*}}{1+\omega z^{*}}e^{m\tau_1}
 +\frac{pe^{m\tau_{1}} y^{*}z(t)}{1+\omega z(t)} \\
&\quad +\frac{pe^{m\tau_1}(z(t)-z^{*})(y(t)z^{*}-y^{*}z(t))}{z(t)(1+\omega z^{*})^{2}}
 -\frac{npy^{*}y(t)e^{m\tau_1}(z(t)-z^{*})^{2}}{bz(t)(1+\omega z^{*})^{2}} \\
&\quad +\frac{\beta x^{*}v^{*}}{1+\alpha v^{*}}\ln\frac{x(t-\tau_{1})
 v(t-\tau_{1})(1+\alpha v(t))}{x(t)v(t)(1+\alpha v(t-\tau_{1}))}
 +\frac{\beta x^{*}v^{*}}{1+\alpha v^{*}}ln\frac{y(t-\tau_{2})}{y(t)} \\
&= -\frac{(x^{*}-x(t))^{2}}{x(t)}\Big(d-r+\frac{rx^{*}}{x_{\rm max}}
 +\frac{rx(t)}{x_{\rm max}}\Big)
 +\frac{\beta x^{*}v^{*}}{1+\alpha v^{*}}\Big(1-\frac{x^{*}}{x(t)}
 +\frac{1+\alpha v^{*}}{1+\alpha v(t)}\frac{v(t)}{v^{*}} \\
&\quad - \frac{ x(t-\tau_{1})v(t-\tau_{1})}{1+\alpha 
 v(t-\tau_{1})}\frac{y^{*}}{y(t)}\frac{1+\alpha v^{*}}{x^{*}v^{*}}+1-\frac{v(t)}{v^{*}}
 -\frac{y(t-\tau_{2})}{y^{*}}\frac{v^{*}}{v(t)}+1 \\
&\quad +\ln\frac{x(t-\tau_{1})v(t-\tau_{1})(1+\alpha v(t))}{x(t)v(t)(1+\alpha 
 v(t-\tau_{1}))}+\ln\frac{y(t-\tau_{2})}{y(t)}\Big)
 -\frac{pe^{m\tau_1}(z(t)-z^{*})(y(t)-y^{*})}{(1+\omega z(t))(1+\omega z^{*})} \\
&\quad +\frac{pe^{m\tau_1}(z(t)-z^{*})(y(t)z^{*}-y^{*}z(t))}{z(t)(1+\omega z^{*})^{2}}
 -\frac{npy^{*}y(t)e^{m\tau_1}(z(t)-z^{*})^{2}}{bz(t)(1+\omega z^{*})^{2}} \\
&= -\frac{(x^{*}-x(t))^{2}}{x(t)}\Big(d-r+\frac{rx^{*}}{x_{\rm max}}
 +\frac{rx(t)}{x_{\rm max}}\Big)+\frac{\beta x^{*}v^{*}}{1+\alpha v^{*}}
 \Big(1-\frac{x^{*}}{x(t)}+\ln\frac{x^{*}}{x(t)} \\
&\quad +1-\frac{ x(t-\tau_{1})v(t-\tau_{1})}{1+\alpha v(t-\tau_{1})}
 \frac{y^{*}}{y(t)}\frac{1+\alpha v^{*}}{x^{*}v^{*}}
 +\ln\frac{ x(t-\tau_{1})v(t-\tau_{1})}{1+\alpha v(t-\tau_{1})}\frac{y^{*}}{y(t)}
 \frac{1+\alpha v^{*}}{x^{*}v^{*}} \\
&\quad +1-\frac{y(t-\tau_{2})}{y^{*}}\frac{v^{*}}{v(t)}+\ln\frac{y(t-\tau_{2})}{y^{*}}\frac{v^{*}}{v(t)}+\ln\frac{1+\alpha v(t)}{1+\alpha v^{*}}
-\frac{1+\alpha v(t)}{1+\alpha v^{*}}+1 \\
&\quad +\frac{1+\alpha v^{*}}{1+\alpha v(t)}\frac{v(t)}{v^{*}}
 -\frac{v(t)}{v^{*}}+\frac{1+\alpha v(t)}{1+\alpha v^{*}}-1\Big)\\
&\quad -\frac{pe^{m\tau_1}(z(t)-z^{*})^{2}(y(t)+\omega y^{*}yz(t))}{z(t)
 (1+\omega z(t))(1+\omega z^{*})}
 -\frac{npy^{*}y(t)e^{m\tau_1}(z(t)-z^{*})^{2}}{bz(t)(1+\omega z^{*})^{2}} \\
&= -\frac{(x^{*}-x(t))^{2}}{x(t)}\Big(d-r+\frac{rx^{*}}{x_{\rm max}}
 +\frac{rx(t)}{x_{\rm max}}\Big)
 -\frac{\beta x^{*}v^{*}}{1+\alpha v^{*}}\Big(g\Big(\frac{x^{*}}{x(t)}\Big) \\
&\quad +g\Big(\frac{ x(t-\tau_{1})v(t-\tau_{1})}{1+\alpha v(t-\tau_{1})}
 \frac{y^{*}}{y(t)}\frac{1+\alpha v^{*}}{x^{*}v^{*}}\Big)
 +g\Big(\frac{y(t-\tau_{2})}{y^{*}}\frac{v^{*}}{v(t)}\Big)
 +g\Big(\frac{1+\alpha v(t)}{1+\alpha v^{*}}\Big)\\
&\quad +\frac{\alpha(v(t)-v^{*})^{2}}{v^{*}(1+\alpha v(t))(1+\alpha v^{*})} \Big)
 -\frac{pe^{m\tau_1}(z(t)-z^{*})^{2}(y(t)+\omega y^{*}yz(t))}{z(t)(1+\omega z(t))
  (1+\omega z^{*})}\\
&\quad -\frac{npy^{*}y(t)e^{m\tau_1}(z(t)-z^{*})^{2}}{bz(t)(1+\omega z^{*})^{2}}.
\end{align*}
Using the fact $x_{\rm max}(r-d)<rx^{*}$, we have 
$\big(d-r+\frac{rx^{*}}{x_{\rm max}}+\frac{rx(t)}{x_{\rm max}}\big)>0$. 
Thus, it follows that $N'(t)\leq 0$ and
$N'(t)= 0$  holds if and only if  $x(t)=x^{*}$, $y(t)=y^{*}$, $v(t)=v^{*}$, 
$z(t)=z^{*}$.
Furthermore, it can be shown that the largest invariant set of 
$\{(x,y,v,z)\!:  N'(t)=0\}$ is the singleton $E^*$. Therefore, 
By LaSalle's invariance  principle \cite{HV}, we see the infection equilibrium  
$ E^{*}$ is  globally asymptotically stable.
\end{proof}

\begin{thebibliography}{99}
\bibitem[1]{BM}  S. Bonhoeffer, R. M. May, G. M. Shaw, M. A. Nowak;
\emph{Virus dynamics and drug therapy}, Proc. Natl. Acad. Sci. USA, 94 (1997) 6971-6976.
\bibitem[2]{BR}   F. A. T. Boshier, D. B. Reeves, E. R. Duke,  et al.;
\emph{Substantial uneven proliferation of CD4$^+$T cells during recovery from acute 
HIV infection is sufficient to explain the observed expanded clones in the HIV reservoir}.
J. Virus. Erad. 8(4) (2022), 100091.
\bibitem[3]{DX}  Y. Deng, R. Xu;
\emph{Stability analysis of a dynamic model of HIV infection with CTL immune response 
and immune impairment},
Acta Math. Sci. 42(A) (2022) 1592-1600.
\bibitem[4]{DW}  P. van den Driessche, J. Watmough;
\emph{Reproduction numbers and sub-threshold endemic equilibria for compartmental 
models of disease transmission}, Math. Biosci. 180 (2002) 29-48.
\bibitem[5]{HV}  J. K. Hale, S. M. Verduyn Lumel;
\emph{Introduction to Functional Differential Equations}, Springer, New York, 1993.
\bibitem[6]{HZ}  Z. X. Hu, J. J. Zhang, H. Wang, F. Liao;
\emph{Dynamics analysis of a delayed viral infection model with logistic growth 
and immune impairment}, Appl. Math. Model. 38 (2014) 524-534.
\bibitem[7]{LGM}  M. Li, K. Guo, W. Ma;
\emph{Uniform persistence and global attractivity in a delayed virus dynamic
model with apoptosis and both virus-to-cell and cell-to-cell infections}. 
 Mathematics,  10 (2022) 975.
\bibitem[9]{K}  A. Korobeinikov;
\emph{Global properties of basic virus dynamics models}. 
Bull. Math. Biol. 66(4) (2004)  879-883.
\bibitem[10]{LS}  M. Y. Li, H. Shu;
\emph{Joint effects of mitosis and intracellular delay on viral dynamics:
two-parameter bifurcation analysis}, J. Math. Biol. 64 (2012) 1005-1020.
\bibitem[11]{LW2}  M. Y. Li, L. Wang;
\emph{Backward bifurcation in a mathematical model for HIV infection in 
vivo with anti-retroviral treatment}, Nonlinear Anal: Real World Appl.
17 (2014) 147-160.
\bibitem[12]{LX}  J. Z. Lin, R. Xu, X. H. Tian;
\emph{Threshold dynamics of an HIV-1 model with both viral and cellular, 
cell-mediated and humoral immune responses}, 
Math. Biosci. Eng. 16 (2019) 292-319.
\bibitem[14]{NA}  M. A. Nowak, R. M. Anderson, M. C. Boerlijst, S. Bonhoeffer, R. M. May;
\emph{HIV-1 evolution and disease progression}, Science. 274 (1996) 1008-1011.
\bibitem[15]{NB}  M. A. Nowak, C. R. M. Bangham;
\emph{Population dynamics of immune response to persistent viruses}, 
Science. 272 (1996) 74-79.
\bibitem[16]{PC}  X. Pan, Y. Chen, H. Shu;
\emph{Rich dynamics in a delayed HTLV-I infection model:
Stability switch, multiple stable cycles, and torus}, 
J. Math. Anal. Appl. 479 (2) (2019) 2214-2235.
\bibitem[17]{QJ}  K. Qi, D. Q. Jiang, T. Hayat,  A. Alsaedi;
\emph{The stationary distribution and extinction of a double thresholds HTLV-I 
infection model with nonlinear CTL immune response disturbed by white noise}, 
Int. J. Biomath. 12 (2019) 1950058.
\bibitem[18]{RW}  R. R. Regoes, D. Wodarz, M. A. Nowak;
\emph{Virus dynamics: the effect of target cell limitation and immune responses 
on virus evolution},  J. Theoret. Biol. 191 (1998) 451-462.
\bibitem[19]{Sa}  G. P. Samanta;
\emph{Permanence and extinction of a nonautonomous HIV/AIDS epidemic model with 
distributed time delay}, Nonlinear Anal: Real World Appl.
 12 (2011) 1163-1177.
\bibitem[20]{S}  H. L. Smith;
\emph{Monotone dynamical systems: An introduction to the theory of competitive and 
cooperative systems}. Mathematical Surveys and Monographs,
 American Mathematical Society.  Providence, RI, 1995.
\bibitem[21]{SN}  X. Song, A. U. Neumann;
\emph{Global stability and periodic solution of the viral dynamics}, 
J. Math. Anal. Appl. 329 (2007) 281-297.
\bibitem[22]{WL}  L. W. Wang, Z. J. Liu,  Y. Li, D. Xu;
\emph{Complete dynamical analysis for a nonlinear HTLV-I infection model with 
distributed delay, CTL response and immune impairment}, 
Discrete Contin. Dyn. Syst. Ser. B. 25 (2020) 917-933.
\bibitem[23]{WS}  S. L. Wang, X. Y. Song, Z. H. Ge;
\emph{Dynamics analysis of a delayed viral infection model with immune impairment}, 
Appl. Math. Model. 35 (2011) 4877-4885.
\bibitem[24]{WLL}  Y. Wang, M. Lu, J. Liu;
\emph{Global stability of a delayed virus model with latent infection and 
Beddington-DeAngelis infection function}, 
Appl. math. letters, 107 (2020) 106463.
\bibitem[25]{X}  R. Xu;
\emph{Global stability of an HIV-1 infection model with saturation infection and 
intracellular delay}. J. Math. Anal. Appl. 375 (2011)  75-81.
\bibitem[26]{ZZZ}  X. Zhou, L. Zhang, T. Zheng, H. L. Li, Z. Teng;
\emph{Global stability for a delayed HIV reactivation model with latent infection and 
Beddington-DeAngelis incidence},
 Appl. Math. Letters, 117 (2021) 107047.
\end{thebibliography}
\end{document}
